\documentclass{article}
\usepackage{geometry}
\geometry{a4paper,scale=0.8}
\usepackage[utf8]{inputenc}

\usepackage{amsmath, amssymb}
\usepackage{amsthm}
\usepackage{enumitem}

\newtheoremstyle{breakstyle}
    {5pt}
    {10pt}
    {\normalfont}
    {0pt}
    {\bfseries}
    {\newline\indent}
    {0pt}
    {\thmname{#1}\thmnumber{ #2}\thmnote{ \textbf{[#3]}}}
\theoremstyle{breakstyle}
\newtheorem{exercise}{Exercise}[section]
\newtheorem{remark}{Remark}[section]

\let\ker\relax
\DeclareMathOperator{\ker}{Ker}
\DeclareMathOperator{\im}{Im}

\title{Chapter 5 - Integral Dependence and Valuations}
\author{Flandre}
\date{20 Jun 2026}

\begin{document}
\maketitle

\setcounter{section}{5}
\begin{exercise}
    Use \emph{Theorem 5.10}, trivial.
\end{exercise}

\begin{exercise}
    Use \emph{Theorem 5.10}, trivial.
\end{exercise}

\begin{exercise}
    Because
    $$
    \im\left( f\otimes 1 \right) =\left( \im{f} \right)\otimes _{A}C 
    $$
    Therefore $B'\otimes _{A}C$ is integral over $\im{f}\otimes _{A}f$, as $B'$ is integral over $\im{f}$.
\end{exercise}

\begin{remark}
    Indeed, \emph{Exercise 5.3} includes \emph{Proposition 5.6(ii)} as special case, as $S^{-1}B\cong S^{-1}A\otimes _{A}B$.
\end{remark}

\begin{exercise}
    Following from the hinted counterexample, the answer is no.
\end{exercise}

\begin{exercise}
    \leavevmode \vspace{-\baselineskip}
    \begin{enumerate}[label=({\roman*})]
        \item 
            Since $x^{-1}\in B,\exists a_{n-1},\ldots ,a_0,\text{s.t. }$
            $$
            x^{-n}+a_{n-1}x^{-n+1}+\ldots +a_{1}x^{-1}+a_0=0
            .$$
            Hence
            $$
            x^{-1}=-a_{n-1}-a_{n-2}x-\ldots -a_1^{n-2}-a_0x^{n-1}\in A
            $$
        \item 
            Use \emph{Corollary 5.8}, trivial.
    \end{enumerate}
\end{exercise}

\begin{exercise}
    Indeed.
    Denote the integral ring homomorphism
    $$
    f_{k}:A\longrightarrow B_{k}
    $$
    and we have
    $$
    \prod_{k=1}^{n} f_{k}:A^{n}\longrightarrow \prod_{k=1}^{n} B_{k}
    $$
    as an integral ring homomorphism.

    (Elimination of terms to make the elements integral is easy, as we can nest those polynomials by terms.)
\end{exercise}

\begin{exercise}
    
\end{exercise}

\end{document}
